%%%%%%%%%%%%%%%%%%%%%%%%%%%%%%%%%%%%%%%%%
% Medium Length Graduate Curriculum Vitae
% LaTeX Template
% Version 1.1 (9/12/12)
%
% This template has been downloaded from:
% http://www.LaTeXTemplates.com
%
% Original author:
% Rensselaer Polytechnic Institute (http://www.rpi.edu/dept/arc/training/latex/resumes/)
%
% Important note:
% This template requires the res.cls file to be in the same directory as the
% .tex file. The res.cls file provides the resume style used for structuring the
% document.
%
%%%%%%%%%%%%%%%%%%%%%%%%%%%%%%%%%%%%%%%%%

%----------------------------------------------------------------------------------------
%	PACKAGES AND OTHER DOCUMENT CONFIGURATIONS
%----------------------------------------------------------------------------------------

\documentclass[margin, 10pt]{res} 
\usepackage{setspace}
\usepackage{hyperref}
\usepackage{helvet} 
\usepackage{graphicx}
\usepackage{enumitem}
\usepackage[english]{babel}
\usepackage[autostyle, english = american]{csquotes}
\usepackage[symbol]{footmisc}

\MakeOuterQuote{"}
\setlength{\textwidth}{5.1in} 

%----------------------------------------------------------------------------------------
%	NAME AND ADDRESS SECTION
%----------------------------------------------------------------------------------------
\begin{document}

\moveleft.5\hoffset\centerline{\large\bf Eghbal A. Hosseini} % Your name at the top
 
\moveleft\hoffset\vbox{\hrule width\resumewidth height 1pt}\smallskip % Horizontal line after name; adjust line thickness by changing the '1pt'

\moveleft.5\hoffset\centerline{ehoseini@mit.edu}
\moveleft.5\hoffset\centerline{(703)789-6512}
 
\moveleft.5\hoffset\centerline{ } % Your address
\moveleft.5\hoffset\centerline{ }
%----------------------------------------------------------------------------------------
\begin{resume}
%----------------------------------------------------------------------------------------
%	EDUCATION SECTION
%----------------------------------------------------------------------------------------
\section{EDUCATION}
\textbf{Ph.D, Computational Neuroscience}   \hfill 2016-2024 \\
Massachusetts Institute of Technology (MIT), Cambridge, MA
\vspace{1mm} \\
Thesis: Towards Synergistic Understanding of Language Processing in Biological and Artificial Systems \\
GPA 5.0/5.0\ 
\vspace{1mm} \\
 Relevant courseworks: Computational Neuroscience (Harvard - MCB131)  Quantitative Methods in Neuroscience (MIT - 9.014), Computational Cognitive Science (MIT - 9.660), Cognitive Science (MIT - 9.012),  How to Make Almost Anything (MIT - MAS 863) , Matrix Methods (MIT - 18.0651)
 
 \textbf{Analytical Connectionism} \hfill 2023\\
Gatsby Computational Neuroscience Unit,   UCL,\
London, UK


\textbf{Brains, Minds, and Machines} \hfill 2017\\
Marine Biological Laboratory, Woods Hole, MA   \
\vspace{2.5mm} \\
\textbf{ MS., Electrical Engineering}  \hfill 2012-2014 \\
George Mason University (GMU), Fairfax, VA \
\vspace{1mm} \\
Thesis: Multi-rate state-dependent primitives underlie the motor adaptation and unlearning to motion dependent force perturbation. \\
GPA: 3.8/4 
\vspace{2.5mm} \\
\textbf{ BS.,  Electrical Engineering}  \hfill 2005-2010 \\
Iran University of Science and Technology (IUST), Tehran, Iran \
\vspace{1mm} \\
Thesis : Position control of DC motor using wavelet based multiresolution analysis. \\
GPA: 16.43/20 (3.41/4)
%----------------------------------------------------------------------------------------
%	HONORS SECTION
%----------------------------------------------------------------------------------------
\section{HONORS \& \\  AWARDS} 
Friends of the McGovern Institute Fellowship, MIT \hfill  2020 \\
BCS Hilibrand Graduate Student Fellowship, MIT \hfill 2017-2018 \\
Henry E. Singleton(1940) Presidential Fellowship, MIT \hfill 2016-2017 \\
ECE Chairmans Award, Volgenau School of engineering, GMU  \hfill Spring 2014 \\
Volgenau School of Engineering Dean Fellowship, GMU  \hfill Spring-Fall 2012 \\
Honors in ECE control group class of 2005, IUST  \hfill Fall 2010 \\
Honors student in ECE class of 2005, IUST  \hfill 2005 \& 2007 
%----------------------------------------------------------------------------------------
%	Patents and publication SECTION
%----------------------------------------------------------------------------------------
\section{PUBLICATIONS}
\textbf{Selected}\ 
\vspace{2mm} \\
\textbf{Eghbal A Hosseini}, Brian Cheung, Evelina Fedorenko, Alex H Williams. 2026. ``Modulating Cross-Modal Convergence with Single-Stimulus, Intra-Modal Dispersion.'' \href{https://openreview.net/forum?id=Ab3pRobOhv}{ICLR 2026 Re-Align Workshop} (accepted). \vspace{2.5mm} \\
\textbf{Eghbal A Hosseini}, Yuxuan Li, Yasaman Bahri, Dylan Campbell, Andrew Kyle Lampinen. 2026. ``Context Structure Reshapes the Representational Geometry of Language Models.'' \href{https://arxiv.org/abs/2601.22364}{arXiv:2601.22364}. \vspace{2.5mm} \\
Andrew Kyle Lampinen, Yuxuan Li, {\bf Eghbal Hosseini}, Sangnie Bhardwaj, Murray Shanahan. 2026. ``Linear representations in language models can change dramatically over a conversation.'' arXiv. \vspace{2.5mm} \\
\textbf{Eghbal A Hosseini}, Colton Casto, Noga Zaslavsky, Colin Conwell, Martin Richardson, Evelina Fedorenko. 2024. ``Universality of representation in biological and artificial neural networks.'' \href{https://www.biorxiv.org/content/10.1101/2024.12.26.629294v1.full}{bioRxiv}. \vspace{2.5mm} \\
\textbf{Eghbal A Hosseini}, Evelina Fedorenko "Large language models implicitly learn to straighten neural sentence trajectories to construct a predictive representation of natural language.", accepted in Neural Information Processing Conference (NeurIPS) , 2023, New Orleans, U.S.A. \vspace{2.5mm} \\
\textbf{Eghbal A Hosseini}, Noga Zaslavsky, Colton Casto, Evelina Fedorenko, "Teasing apart the representational spaces of ANN language models to discover key axes of model-to-brain alignment" {Contributed talk - top 5\% of submission, Computational Cognitive Neuroscience (CCN), 2023, Oxford, England }.  \vspace{2.5mm} \\
\textbf{Eghbal A Hosseini}, Martin A Schrimpf, Yian Zhang, Samuel Bowman, Noga Zaslavsky, Evelina Fedorenko. 2022. “Artificial Neural Network Language Models Predict Human Brain Responses to Language Even After a Developmentally Realistic Amount of Training.” . \href{https://doi.org/10.1162/nol\_a\_00137}{Neurobiology of Language} 2024; 5 (1): 43–63.  \vspace{2.5mm} \\
Schrimpf, Martin, Idan A. Blank, Greta Tuckute, Carina Kauf, {\bf Eghbal A. Hosseini}, Nancy G. Kanwisher, Joshua B. Tenenbaum, and Evelina Fedorenko. 2021. “The neural architecture of language: Integrative reverse-engineering converges on a model for predictive processing” . \href{https://doi.org/10.1101/2020.06.26.174482.}{PNAS}  \vspace{2.5mm} \\
\textbf{Continued} 
\vspace{2mm} \\
Tamar I Regev$^\ast$, Colton Casto$^\ast$, \textbf{Eghbal A Hosseini}, Markus Adamek, Peter Brunner, Evelina Fedorenko. 2022. “Intracranial recordings reveal three distinct neural response patterns in the language network” . \href{https://www.biorxiv.org/content/10.1101/2022.12.30.522216v1.abstract}{BioRxiv}  \vspace{2.5mm} \\
Wang, Jing, {\bf Eghbal Hosseini }, Nicolas Meirhaeghe, Adam Akkad, and Mehrdad Jazayeri. 2020. “Reinforcement Regulates Timing Variability in Thalamus.” eLife 9 (December). https://doi.org/10.7554/eLife.55872. \vspace{2.5mm} \\
Tremblay, Sébastien, Leah Acker, Arash Afraz, Daniel L. Albaugh, Hidetoshi Amita, Ariana R. Andrei, Alessandra Angelucci,...,{\bf Eghbal A. Hosseini},... et al. 2020. “An Open Resource for Non-Human Primate Optogenetics.” Neuron, \vspace{2.5mm} \\
Alhussein, Laith, {\bf Eghbal A. Hosseini }, Katrina P. Nguyen, Maurice A. Smith, and Wilsaan M. Joiner. 2019. “Dissociating Effects of Error Size, Training Duration, and Amount of Adaptation on the Ability to Retain Motor Memories.” Journal of Neurophysiology 122 (5): 2027–42. \vspace{2.5mm} \\
Nguyen K.P, Weiwei Z., McKenna E. L., Colucci K.,, Bray L. C.,{ \bf Hosseini E.A.}, Alhussein L., Joiner W. M., 2019 "The 24 hour savings of adaptation to novel movement dynamics initially reflects the recall of previous performance",  Journal of Neurophysiology,  \ 
\vspace{2.5mm} \\
Wang, Jing$^\ast$ Devika Narain\footnote{co-first authors}, {\bf Eghbal A. Hosseini}, and Mehrdad Jazayeri. 2018. "Flexible Timing by Temporal Scaling of Cortical Responses." \href{https://www.nature.com/articles/s41593-017-0028-6}{ Nature Neuroscience 21 (1):102-10}.
\vspace{2.5mm} \\
Remington, Evan D., Devika Narain, {\bf Eghbal A. Hosseini }, and Mehrdad Jazayeri. 2018. "Flexible Sensorimotor Computations through Rapid Reconfiguration of Cortical Dynamics." \href{https://www.cell.com/neuron/abstract/S0896-6273(18)30418-5}{ Neuron 98 (5). Elsevier: 1005-19.e5.}  \ 
\vspace{2.5mm} \\
\textbf{Eghbal A Hosseini}, Katrina P. Nguyen, and Wilsaan M. Joiner. 2017. "The Decay of Motor Adaptation to Novel Movement Dynamics Reveals an Asymmetry in the Stability of Motion State-Dependent Learning." \href{https://doi.org/10.1371/journal.pcbi.1005492} { PLoS Computational Biology 13 (5): e1005492.} \
\vspace{2.5mm} \\
\textbf{Eghbal A Hosseini}, and H. Sadjadian. 2015. "Noise Resistant Design of Wavelet Based Multiresolution Control."  \href{http://ieeexplore.ieee.org/document/7172111/} {In American Control Conference (ACC), 2015, 4959?63.}  \
\vspace{5mm} \\
\textbf{Posters - Presentations} \
\vspace{2mm} \\
\textbf{Eghbal A Hosseini}, Martin A Schrimpf, Yian Zhang, Samuel Bowman, Noga Zaslavsky, Evelina Fedorenko. 2022. “Alignment of ANN Language Models with Humans After a Developmentally Realistic Amount of Training” .Cosyne 2023 ,Montreal, Canada \vspace{2.5mm} \\
 {\bf  Hosseini E.A}, Schrimpf M., Bowman S., Fedorenko E., Zaslavsky N. "The effect of training in neural network language models on predicting brain activity" Society for Neurobiology of Language,  2020. \vspace{2.5mm} \\
 Wang J. , {\bf  Hosseini E.A},  Meirhaeghe N., Akkad A., and  Jazayeri M. , "Reinforcement regulates context-dependent timing variability in thalamus", Cosyne 2020 , Denver, CO \vspace{2.5mm} \\
Remington E. D., Narain D. { \bf Hosseini E.A.}, Jazayeri M., "Control of sensorimotor dynamics through adjustment of inputs and initial condition",
 Cosyne 2018 , Denver, CO \vspace{2.5mm} \\
 Wang J., { \bf Hosseini E.A.}, Jazayeri M., "Reward-dependent modulation of variability mediates trial-by-trial motor learning",
Society for Neuroscience meeting, 2018, San Diego, CA. \
\vspace{2.5mm} \\
Wang J., Jazayeri M., { \bf Hosseini E.A.}, Narain D., "The speed of neural dynamics as a neural code for motor timing", Computational and System Neuroscience Meeting (Cosyne), 2017, Salt Lake City, UT \
\vspace{2.5mm} \\
{ \bf Hosseini E.A.}, Wang J., Jazayeri M., "Representation of contextual information in cortico-basal ganglia circuits during motor timing",
Society for Neuroscience meeting, 2016, San Diego, CA. \
\vspace{2.5mm} \\
Wang J., { \bf Hosseini E.A.}, Jazayeri M., "Scalar dynamics in neural activity during timing",
Society for Neuroscience meeting, 2016, San Diego, CA. \
\vspace{2.5mm} \\
Remington E. D., { \bf Hosseini E.A.}, Jazayeri M., "Probing a sensorimotor transformation in dorsomedial frontal cortex using electrophysiology and optogenetics",
Society for Neuroscience meeting, 2016, San Diego, CA. \
\vspace{2.5mm} \\
{ \bf Hosseini E.A.}, Nguyen K.P., Joiner W.M., "Multi-rate state-dependent primitives underlie the motor adaptation and unlearning to motion dependent force perturbation", McGovern Institute Spring Symposium, MIT, 2015, Cambridge, MA. \
\vspace{2.5mm} \\
Remington E. D., { \bf Hosseini E.A.}, Jazayeri M., "Sensory measurement and motor planning are not 
separable in interval timing", Society for Neuroscience, 2015, Chicago, IL.\
\vspace{2.5mm} \\
Nguyen K.P, McKenna E. L., Bray L. C., Colucci K., Alhussein L. { \bf Hosseini E.A.},  Joiner W. M., "The initial single-trial rate of motor adaptation savings is dependent on
both the training duration and final adaptive state before a 24-hour break", Society for Neuroscience, 2015, Chicago, IL.\
\vspace{2.5mm} \\
Alhussein L., { \bf Hosseini E.A.}, Nguyen K.P., Joiner W.M., "The Intralimb stability
of adaptation to novel movement dynamics is dictated by the training duration for different types of motion-dependent perturbations", Neural Control of Movement Conference, 2015, Charleston, SC. \
\vspace{2.5mm} \\
Nguyen K.P, { \bf Hosseini E.A.}, Joiner W.M., "The decay of motor adaptation to novel
movement dynamics reveals hysteresis in motor primitive gain-space", Society for Neuroscience, 2014, Washington D.C. \
\vspace{2.5mm} \\
Keshtkar H., Sartipizadeh H.,{ \bf Hosseini E.A.}, Khandani A., Naghavi F., "The role of telework centers in development of electronic municipality", 1st international conference on electronic municipality, 2007, Tehran, Iran. \
\vspace{2.5mm} \\
Keshtkar H.,{ \bf Hosseini E.A.}, "Telework centers and economic productivity",  National conference on industry, student and sustainable improvement, 2007, Tehran,
Iran. \
\vspace{5mm} \\
 \textbf{Patents}  
 \vspace{2.5mm} \\
{ \bf Hosseini E.A.}, Momtazan H., Momtazan A., "Automatic device for electric arc based production of carbon nanotubes in Liquid environment" Patent Number 72901, 2011, Iran. \

\section { RESEARCH  \\ EXPERIENCE} 

\textbf{Visiting Researcher}  \hfill 2024-Present\\ 
Google DeepMind, New York, NY
\vspace{2.5mm} \\
\textbf{Graduate Research Assistant}  \hfill Summer 2020-2024\\ 
Dr. Evelina Fedorenko, EvLab, McGovern Institute for Brain Research, MIT 
\vspace{2.5mm} \\
\textbf{Graduate Fellowship Student}  \hfill Spring 2019-Fall 2019\\ 
Dr. Evelina Fedorenko, EvLab, McGovern Institute for Brain Research, MIT 
\vspace{2.5mm} \\
\textbf{Graduate Fellowship Student}  \hfill Spring 2019\\ 
Dr. Ila Fiete, Fiete Lab, McGovern Institute for Brain Research, MIT 
\vspace{2.5mm} \\
\textbf{Graduate Fellowship Student}  \hfill Summer 2017-Summer 2018\\ 
Dr. Edward S. Boyden, Synthetic Neurbiology Group, McGovern Institute for Brain Research, MIT 
\vspace{2.5mm} \\
\textbf{Technical Assistant}  \hfill Spring 2015-Summer 2016\\ 
Dr. Mehrdad Jazayeri, JazLab, McGovern Institute for Brain Research, MIT 
\vspace{2.5mm} \\
\textbf{Graduate Research Assistant}  \hfill Spring 2013-Fall 2014\\ 
Dr. Wilsaan Joiner, Sensorimotor Integration Lab, Volgenau School of Engineering, GMU 
\vspace{2.5mm} \\
\textbf{Undergraduate Research Assistant}  \hfill 2009-2010\\ 
Mechatronic and Robotic Research Lab, IUST 
\vspace{2.5mm} \\
\textbf{Undergraduate Research Assistant}  \hfill 2006-2007\\ 
Electronics Research Center, IUST 
\section { TEACHING  \\ EXPERIENCE} 
\textbf{Teaching Assistant}  \hfill Sprint 2020\\
Introduction to Neural Computation, Department of Brain and Cognitive Sciences, MIT \
\vspace{2.5mm} \\
\textbf{Teaching Assistant}  \hfill Fall 2017\\
Science of Intelligence, Department of Brain and Cognitive Sciences, MIT \
\vspace{2.5mm} \\
\textbf{Teaching Assistant}   \hfill Fall 2014\\
Introduction to Biomedical Engineering, Bioengineering department, GMU \
\vspace{2.5mm} \\
\textbf{Graduate Research Assistant}  \hfill Summer 2012\\ 
Center for Outreach in Mathematics Professional Learning and Educational \\
Technology (COMPLETE), GMU
\vspace{2mm} \
\begin{itemize} \itemsep -2pt % Reduce space between items
\item  Designed a series of experiments for demonstrating the use of high school physics and calculus in solving engineering problems, and mentored high school teachers as
they implemented these experiments in the coursework of two 10th grade classes in Northern Virginia high schools. 
\end{itemize}
\textbf{Teaching Assistant}   \hfill Spring-Fall 2012\\
Bioengineering Measurements Lab, Bioengineering department, GMU \
\vspace{2.5mm} \\
\textbf{Teaching Assistant}  \hfill Spring 2008- Fall 2010\\ 
Circuit Theory, Department of Computer Engineering, IUST   
%---------------------------------------------------------------------------------------
%	PROFESSIONAL EXPERIENCE SECTION
%----------------------------------------------------------------------------------------
\section{COMPUTER \\ SKILLS}
\textbf{Languages \& Software: } MATLAB, Simulink, Python, Tensorflow,Pytorch, R, PSPICE, Protel DXP, Microsoft
Office, Adobe Illustrator, Adobe Acrobat Pro. Solidworks, LaTex, Slurm
\vspace{2.5mm} \\
\textbf{Operating Systems: } Linux (Ubuntu), Macintosh OS, Microsoft Windows
\section{LANGUAGES}
English (fluent) \\
Farsi (native)\
\section { PROFESSIONAL  \\ MEMBERSHIP} 
Society for Neuroscience \hfill 2013-2019\\ 
Institute of Electrical and Electronic Engineers (IEEE)  \hfill 2015-2018\\ 
\end{resume}
\end{document}